\documentclass[11pt]{article}
\usepackage[margin=1in]{geometry}
\usepackage{amsmath,amssymb,amsthm,mathtools}
\usepackage{graphicx}
\usepackage{microtype}
\usepackage[colorlinks=true,linkcolor=blue,citecolor=blue,urlcolor=blue]{hyperref}

\newtheorem{theorem}{Theorem}
\newtheorem{lemma}[theorem]{Lemma}
\newtheorem{corollary}[theorem]{Corollary}
\newcommand{\C}{\mathbb C}
\newcommand{\Span}{\operatorname{span}}
\newcommand{\rank}{\operatorname{rank}}
\newcommand{\udim}{\overline{\dim}_{B}}

\title{Finite Box Dimension Forces Finite-Rank\\Taylor Coefficients}
\author{Automated mathematics run \texttt{fa\_banach\_001}, lane 10}
\date{13 August 2026}

\begin{document}
\maketitle

\begin{center}
\fbox{\begin{minipage}{0.88\textwidth}
\textbf{Status.} Candidate full affirmative solution to Questions 4.4 and
4.5 of Carando--D'Andrea--Torres--Turco, likely valid, pending human review.
The proof gives the stronger conclusion that every homogeneous Taylor
coefficient has finite-dimensional range, with an explicit rank bound.
\end{minipage}}
\end{center}

\section{Source questions and answer}

Let $E,F$ be complex Banach spaces, let $U\subset E$ be open, and let
$f:U\to F$ be holomorphic.  Write
\[
 f(x_0+x)=\sum_{m=0}^{\infty}P_m(x)
\]
for its Taylor expansion at $x_0$, where $P_m=P_mf(x_0)$ is
$m$-homogeneous.  Questions 4.4 and 4.5 of \cite{CDTT} ask whether
\[
 \dim_B f(x_0+\varepsilon B_E)<\infty
\]
forces, for every $m$, respectively
\[
 (e_n(P_m(B_E)))_{n\ge1}\in\ell_1
 \quad\text{and}\quad
 \udim P_m(B_E)<\infty.
\]
The printed questions write $f:U\to E$ although $F$ has just been introduced;
as everywhere else in the paper, the well-typed formulation is $f:U\to F$.

The answer to both questions is \emph{yes}.  In fact, if
\[
 d:=\udim f(x_0+\varepsilon B_E)<\infty,
 \qquad r:=\left\lfloor\frac d2\right\rfloor,
\]
and
\[
 V_m:=\Span\bigl(P_1(E)\cup\cdots\cup P_m(E)\bigr),
 \qquad V_0=\{0\},
\]
then
\begin{equation}\label{quantitative}
 \dim(V_m/V_{m-1})\le (r+m)^m,
 \qquad
 \dim V_m\le\sum_{k=1}^m(r+k)^k.
\end{equation}

\clearpage
\begin{figure}[p]
\centering
\includegraphics[width=\textwidth]{figures/open_problem_crop.png}
\caption{Official arXiv:2401.12059v1 PDF page 15: the end of Proposition 4.2
and Questions 4.4--4.5.}
\end{figure}
\clearpage

\section{Two rank lemmas}

The first observation converts metric dimension into a uniform differential
rank bound.  Complex rank contributes twice its value to real box dimension.

\begin{lemma}[box dimension bounds differential rank]\label{metricrank}
Put $K=f(x_0+\varepsilon B_E)$ and suppose $d=\udim K<\infty$.  Then
\[
 \rank_{\C}Df(y)\le \left\lfloor d/2\right\rfloor
 \qquad(y\in x_0+\varepsilon B_E).
\]
\end{lemma}

\begin{proof}
Suppose $Df(y)$ has rank at least $s$.  Choose $h_1,\ldots,h_s\in E$ and
$\lambda_1,\ldots,\lambda_s\in F^*$ so that the matrix
\[
 \bigl(\lambda_i(Df(y)h_j)\bigr)_{i,j=1}^s
\]
is invertible.  On $H=\Span\{h_1,\ldots,h_s\}$ consider
\[
 \Phi(z)=\bigl(\lambda_1(f(y+z)),\ldots,
                   \lambda_s(f(y+z))\bigr).
\]
After shrinking a neighborhood of $0$ in $H$, it stays inside the given
ball.  The finite-dimensional holomorphic inverse theorem says that its image
contains a nonempty open subset of $\C^s$.  Hence the Lipschitz image
$(\lambda_1,\ldots,\lambda_s)(K)$ has upper box dimension at least $2s$.
Lipschitz maps do not increase upper box dimension, so $2s\le d$.  Taking
$s=\lfloor d/2\rfloor+1$ proves the claim.
\end{proof}

The second lemma is the Banach-space form of the Jacobian estimate proved in
Section 3 of \cite{CDTT}.  We include the reduction because it is exactly what
the Taylor argument needs.

\begin{lemma}[bounded derivative rank forces finite polynomial rank]
\label{polyrank}
Let $P:E\to G$ be a continuous $m$-homogeneous polynomial and suppose
\[
 \rank DP(x)\le r\qquad(x\in E).
\]
Then
\[
 \dim\Span P(E)\le(r+m)^m.
\]
\end{lemma}

\begin{proof}
If $P(x_1),\ldots,P(x_N)$ are linearly independent, choose
$\lambda_1,\ldots,\lambda_N\in G^*$ with
$\lambda_i(P(x_j))=\delta_{ij}$.  Define $L:\C^N\to E$ by
$L(a)=\sum_j a_jx_j$ and define the scalar $m$-homogeneous polynomials
\[
 p_i=\lambda_i\circ P\circ L\qquad(1\le i\le N).
\]
They are linearly independent because $p_i(e_j)=\delta_{ij}$.  The
Jacobian-rank estimate of \cite[Corollary 3.3]{CDTT} gives
\[
 \max_{a\in\C^N}\rank D(p_1,\ldots,p_N)(a)
 \ge N^{1/m}-m.
\]
On the other hand, each of these differentials factors through $DP(L(a))$,
so its rank is at most $r$.  Thus $N\le(r+m)^m$.  If $\Span P(E)$ had larger
dimension, one could choose more than $(r+m)^m$ independent values, a
contradiction.
\end{proof}

For completeness, the estimate invoked above comes from a triangular
coefficient count.  If the maximum Jacobian rank of $N$ linearly independent
$m$-homogeneous scalar polynomials is $q$, choose a nonzero $q$-minor.  Every
polynomial in their span satisfies the associated first-order determinantal
equations.  After a generic linear change of variables, coefficient comparison
in lexicographic order shows that a degree-$m$ solution is determined by the
monomials in the first $q$ variables.  There are
$\binom{q+m-1}{m}\le(q+m)^m$ of these.  Hence
$N\le(q+m)^m$, which is precisely the displayed lower bound.  This also
clarifies that only the final estimate of Corollary 3.3 is being used; no
infinite-dimensional algebraic geometry is hidden in Lemma \ref{polyrank}.

\section{Taylor filtration and the noncancellation determinant}

\begin{theorem}[quantitative affirmative answer]\label{main}
Under the hypotheses in Section 1, assume only that
$d=\udim f(x_0+\varepsilon B_E)<\infty$.  Then every $V_m$ is finite
dimensional and the bounds \eqref{quantitative} hold.  In particular,
each $P_m$ has finite-dimensional range and $P_m(B_E)$ has finite upper box
dimension.
\end{theorem}

\begin{proof}
Set $r=\lfloor d/2\rfloor$.  We argue by induction on $m$.  Suppose
$V_{m-1}$ is finite dimensional.  It is then closed, so the quotient
$G_m=F/V_{m-1}$ is a Banach space.  Let
\[
 \pi_m:F\to G_m,
 \qquad Q_k=\pi_m\circ P_k,
 \qquad g=\pi_m\circ f.
\]
By construction, $Q_k=0$ for $k<m$.

We claim that
\begin{equation}\label{rankQ}
 \rank DQ_m(x)\le r\qquad(x\in E).
\end{equation}
If not, put $s=r+1$ and choose $h_1,\ldots,h_s\in E$ and
$\mu_1,\ldots,\mu_s\in G_m^*$ such that
\[
 M=\bigl(\mu_i(DQ_m(x)h_j)\bigr)_{i,j=1}^s
 \quad\text{has}\quad \det M\ne0.
\]
For sufficiently small $t$, the point $x_0+tx$ lies in the prescribed ball,
and termwise differentiation of the Taylor series gives
\[
 Dg(x_0+tx)h
 =\sum_{k=m}^{\infty}t^{k-1}DQ_k(x)h.
\]
Therefore the analytic scalar function
\[
 \Delta(t)=\det\bigl(\mu_i(Dg(x_0+tx)h_j)\bigr)_{i,j=1}^s
\]
satisfies
\begin{equation}\label{leading}
 \Delta(t)=t^{s(m-1)}\det M+O\bigl(t^{s(m-1)+1}\bigr).
\end{equation}
It is nonzero for all sufficiently small nonzero $t$.  Consequently
$Dg(x_0+tx)$ has rank at least $s$.  Since $Dg=\pi_m\circ Df$, this forces
$Df(x_0+tx)$ to have rank at least $s=r+1$, contradicting Lemma
\ref{metricrank}.  This proves \eqref{rankQ}.

Lemma \ref{polyrank}, applied to $Q_m$, now gives
\[
 \dim\Span Q_m(E)\le(r+m)^m.
\]
But $\Span Q_m(E)$ is canonically $V_m/V_{m-1}$.  This proves the first
bound in \eqref{quantitative}; summing the successive dimensions proves the
second and closes the induction.

Finally $P_m(B_E)$ is a bounded subset of the finite-dimensional complex
space $\Span P_m(E)\subseteq V_m$.  Its upper box dimension is therefore
finite (indeed, at most $2\dim\Span P_m(E)$).
\end{proof}

The quotient in this proof is essential.  Without it, an
$s\times s$ minor of $Df(x_0+tx)$ may contain lower Taylor degrees, and its
coefficient at order $s(m-1)$ can cancel the minor coming from $DP_m(x)$.
Passing to $F/V_{m-1}$ removes every lower degree and makes the leading
coefficient in \eqref{leading} exact.

\begin{corollary}[the missing $\ell_1$ endpoint]\label{entropy}
For every $m\ge1$,
\[
 \bigl(e_n(P_m(B_E))\bigr)_{n\ge1}\in\ell_1.
\]
Thus Questions 4.4 and 4.5 both have affirmative answers.
\end{corollary}

\begin{proof}
Let $D_m=\dim\Span P_m(E)<\infty$.  If $D_m=0$ there is nothing to prove.
Otherwise the standard volumetric covering estimate in the real
$2D_m$-dimensional normed space gives, for a constant $C_m$,
\[
 e_n(P_m(B_E))\le C_m,2^{-n/(2D_m)}.
\]
The right-hand side is summable.
\end{proof}

\section{Verification, novelty, and review focus}

The exact checker bundled with this packet verifies, over the integers, the
leading exponent and leading determinant in \eqref{leading} for matrix sizes
$1$ through $4$ and Taylor degrees $1$ through $5$.  It also checks the
cumulative bounds for representative box dimensions.  These computations
guard the exponent and indexing; the general proof is the determinant
argument above.

Six focused attempts are recorded in the run note.  Direct entropy and convex
hull methods reproduce the source's $p>1$ loss; diagonal counterexample
families cannot place infinite rank in one fixed homogeneous degree.  The
successful route instead uses differential rank, lower-degree quotients, and
the source's algebraic Jacobian estimate.

A bounded literature search on 13 August 2026 covered the exact title,
Questions 4.4--4.5, arXiv id, DOI, authors, and combinations of ``finite box
dimension'', ``Taylor coefficient'', ``finite rank'', and ``entropy
endpoint''.  It found the 2025 Mathematische Nachrichten publication and one
indexed citing paper on generalized entropy numbers, but no claimed answer to
either question.  Novelty confidence is moderate because citation databases
and web search are not exhaustive.

Human review should prioritize Lemma \ref{polyrank}'s reduction to Corollary
3.3 of the source and the quotient determinant in \eqref{leading}.  The
inverse-function and finite-dimensional entropy steps are standard.  The
argument answers a slightly stronger hypothesis than printed: finite
\emph{upper} box dimension suffices, whether or not the box-dimension limit
exists.

\begin{thebibliography}{9}
\bibitem{CDTT}
D. Carando, C. D'Andrea, L. A. Torres, and P. Turco,
\emph{Entropy numbers and box dimension of polynomials and holomorphic
functions}, Mathematische Nachrichten 298 (2025), 567--580;
arXiv:2401.12059v1.
\end{thebibliography}

\end{document}
